% ====================== 主文件导言区必须加载的宏包 ======================
% \usepackage{amsmath,amssymb}
% \usepackage{tikz}
% \usepackage{circuitikz}
% \usepackage{caption} % 用于 \captionof 表格小标题
% 本文件为章节子文件，宏包请在主.tex中统一加载，无需重复导入
% ======================================================================

\chapter{逻辑代数基础}
\setlength{\parskip}{0.4em}
\setlength{\itemsep}{0.2em}
\setlength{\tabcolsep}{4pt} % 统一表格列间距

\section{概述}
\subsection{二值逻辑与逻辑运算}
\begin{itemize}
    \item \textbf{二值逻辑}：采用0、1代表两种对立逻辑状态，仅存在两种取值的逻辑体系。
    \item \textbf{逻辑运算}：对0、1逻辑变量按因果关系进行运算。
\end{itemize}

\subsection{数字电路的特点与描述工具}
\begin{itemize}
    \item \textbf{数字电路}：输入输出仅高低电平，可用0/1二值变量描述，输入输出满足逻辑函数关系。
    \item \textbf{逻辑代数（开关代数）}：布尔代数适配数字电路的形式化工具，是逻辑函数分析基础。
\end{itemize}

\section{逻辑代数中的三种基本运算}
基础逻辑运算：与（AND）、或（OR）、非（NOT）。

\subsection{三种基本逻辑运算的定义}
\subsubsection{与运算（逻辑乘）}
全部输入为1，输出才为1，缺一不可。

\subsubsection{或运算（逻辑加）}
任意输入为1，输出即为1，有一即可。

\subsubsection{非运算（逻辑求反）}
输入取相反状态。

\begin{figure}[!htbp]
\centering
\begin{tikzpicture}[scale=0.75,line width=0.8pt]
% (a) 与门定义电路
\draw (0,0)--(0,1.3);
\draw (-0.15,1.3)--(0.15,1.3);
\draw (-0.3,1.6)--(0.3,1.6);
\draw (0,1.6)--(0,3);
\draw (0,3)--(1.5,3);
\draw (1.5,3)--(1.9,3);
\draw (1.9,3)--(2.45,3.35);
\draw (2.6,3)--(3.0,3);
\node at (2.25,3.7) {$A$};
\draw (3.0,3)--(3.5,3);
\draw (3.5,3)--(4.05,3.35);
\draw (4.2,3)--(4.6,3);
\node at (3.85,3.7) {$B$};
\draw (4.6,3)--(5.3,3);
\draw (5.3,2) circle(0.45);
\draw (5.0,1.7)--(5.6,2.3);
\draw (5.0,2.3)--(5.6,1.7);
\node[right] at (5.8,2) {$Y$};
\draw (5.3,3)--(5.3,2.45);
\draw (5.3,1.55)--(5.3,0);
\draw (5.3,0)--(0,0);
\node at (2.7,-0.9) {(a)};
\node at (2.7,-1.5) {与门定义电路};

% (b) 或门定义电路
\begin{scope}[xshift=8cm]
\draw (0,0)--(0,1.3);
\draw (-0.15,1.3)--(0.15,1.3);
\draw (-0.3,1.6)--(0.3,1.6);
\draw (0,1.6)--(0,3);
\draw (0,3)--(1.2,3);
\fill (1.2,3) circle(0.06);
\draw (1.2,3)--(1.2,4.0);
\draw (1.2,4.0)--(2.0,4.0);
\draw (2.0,4.0)--(2.55,4.35);
\draw (2.7,4.0)--(3.5,4.0);
\node at (2.35,4.7) {$A$};
\draw (1.2,3)--(1.2,2.0);
\draw (1.2,2.0)--(2.0,2.0);
\draw (2.0,2.0)--(2.55,2.35);
\draw (2.7,2.0)--(3.5,2.0);
\node at (2.35,2.7) {$B$};
\draw (3.5,4.0)--(3.5,3);
\draw (3.5,2.0)--(3.5,3);
\fill (3.5,3) circle(0.06);
\draw (3.5,3)--(5.0,3);
\draw (5.0,2) circle(0.45);
\draw (4.7,1.7)--(5.3,2.3);
\draw (4.7,2.3)--(5.3,1.7);
\node[right] at (5.5,2) {$Y$};
\draw (5.0,3)--(5.0,2.45);
\draw (5.0,1.55)--(5.0,0);
\draw (5.0,0)--(0,0);
\node at (2.5,-0.9) {(b)};
\node at (2.5,-1.5) {或门定义电路};
\end{scope}

% (c) 非门定义电路
\begin{scope}[xshift=16cm]
\draw (0,0)--(0,1.3);
\draw (-0.15,1.3)--(0.15,1.3);
\draw (-0.3,1.6)--(0.3,1.6);
\draw (0,1.6)--(0,3);
\draw (0,3)--(1.5,3);
\draw (1.5,2.8) rectangle (2.4,3.2);
\node at (1.95,2.35) {$R$};
\draw (2.4,3)--(4.2,3);
\fill (4.2,3) circle(0.06);
\draw (4.2,3)--(5.3,3);
\draw (5.3,2) circle(0.45);
\draw (5.0,1.7)--(5.6,2.3);
\draw (5.0,2.3)--(5.6,1.7);
\node[right] at (5.8,2) {$Y$};
\draw (5.3,3)--(5.3,2.45);
\draw (5.3,1.55)--(5.3,0);
\draw (5.3,0)--(0,0);
\fill (4.2,0) circle(0.06);
\draw (4.2,3)--(4.2,1.5);
\draw (4.2,1.5)--(4.55,0.9);
\node[left] at (4.0,1.2) {$A$};
\draw (4.2,0)--(4.2,0.8);
\end{scope}
\node at (18.6,-0.9) {(c)};
\node at (18.6,-1.5) {非门定义电路};
\end{tikzpicture}
\caption{基本逻辑运算的定义电路}
\label{fig:logic_definition}
\end{figure}

\subsection{真值表、表达式及逻辑符号}
\subsubsection{与逻辑}
\begin{itemize}
    \item 表达式：$Y = A \cdot B = AB$
    \item 规则：全1出1，有0出0
\end{itemize}

\subsubsection{或逻辑}
\begin{itemize}
    \item 表达式：$Y = A + B$
    \item 规则：有1出1，全0出0
\end{itemize}

\subsubsection{非逻辑}
\begin{itemize}
    \item 表达式：$Y = \overline{A} = A'$
    \item 规则：入0出1，入1出0
\end{itemize}

\begin{figure}[!htbp]
    \centering
    \begin{minipage}{0.3\textwidth}
        \centering
        \captionof{table}{与逻辑真值表}
        \begin{tabular}{|c|c|c|}
            \hline
            $A$ & $B$ & $Y$ \\
            \hline
            0 & 0 & 0 \\
            0 & 1 & 0 \\
            1 & 0 & 0 \\
            1 & 1 & 1 \\
            \hline
        \end{tabular}
    \end{minipage}
    \hfill
    \begin{minipage}{0.3\textwidth}
        \centering
        \captionof{table}{或逻辑真值表}
        \begin{tabular}{|c|c|c|}
            \hline
            $A$ & $B$ & $Y$ \\
            \hline
            0 & 0 & 0 \\
            0 & 1 & 1 \\
            1 & 0 & 1 \\
            1 & 1 & 1 \\
            \hline
        \end{tabular}
    \end{minipage}
    \hfill
    \begin{minipage}{0.3\textwidth}
        \centering
        \captionof{table}{非逻辑真值表}
        \begin{tabular}{|c|c|}
            \hline
            $A$ & $Y$ \\
            \hline
            0 & 1 \\
            1 & 0 \\
            \hline
        \end{tabular}
    \end{minipage}
\end{figure}

\subsubsection{三种基本运算的逻辑符号}
两类标准符号：国标矩形符号、ANSI异形符号。

\begin{figure}[!htbp]
\centering
\begin{tikzpicture}[scale=1]
% 与门
\draw (0,0) rectangle (2,2);
\node at (1,1) {\Large \&};
\draw (-0.8,1.5)--(0,1.5);
\draw (-0.8,0.5)--(0,0.5);
\draw (2,1)--(2.8,1);
\node[left] at (-0.8,1.5) {$A$};
\node[left] at (-0.8,0.5) {$B$};
\node[right] at (2.8,1) {$Y$};
\node at (1,-0.7) {与门};
% 或门
\begin{scope}[xshift=5cm]
\draw (0,0) rectangle (2,2);
\node at (1,1) {$\geq1$};
\draw (-0.8,1.5)--(0,1.5);
\draw (-0.8,0.5)--(0,0.5);
\draw (2,1)--(2.8,1);
\node[left] at (-0.8,1.5) {$A$};
\node[left] at (-0.8,0.5) {$B$};
\node[right] at (2.8,1) {$Y$};
\node at (1,-0.7) {或门};
\end{scope}
% 非门
\begin{scope}[xshift=10cm]
\draw (0,0) rectangle (2,2);
\node at (1,1) {$1$};
\draw (-0.8,1)--(0,1);
\draw (2,1)--(2.0,1);
\draw (2.15,1) circle (0.12);
\draw (2.27,1)--(2.8,1);
\node[left] at (-0.8,1) {$A$};
\node[right] at (2.8,1) {$Y$};
\node at (1,-0.7) {非门};
\end{scope}
\node[font=\bfseries] at (6,-1.8) {（a）国标矩形符号};
\end{tikzpicture}

\vspace{1cm}
\begin{circuitikz}
\draw (0,0) node[and port] (and1) {}
(and1.in 1) -- ++(-1,0)
(and1.in 2) -- ++(-1,0)
(and1.out) -- ++(1,0);
\node at (0,-1.5) {与门};
\draw (5,0) node[or port] (or1) {}
(or1.in 1) -- ++(-1,0)
(or1.in 2) -- ++(-1,0)
(or1.out) -- ++(1,0);
\node at (5,-1.5) {或门};
\draw (10,0) node[not port] (not1) {}
(not1.in) -- ++(-1,0)
(not1.out) -- ++(1,0);
\node at (10,-1.5) {非门};
\node[font=\bfseries] at (5,-2.6) {（b）ANSI异形符号};
\end{circuitikz}
\caption{基础逻辑门两种标准符号}
\label{fig:logic_symbol_compare}
\end{figure}

% ====================== 第四节 重构优化完整部分 ======================
\section{常用复合逻辑运算}
本节介绍五类工程常用复合逻辑：与非、或非、与或非、异或、同或，分别给出逻辑表达式、运算规则、真值表与标准逻辑符号。

\subsection{五类复合逻辑运算定义与运算规则}
\subsubsection{与非运算（NAND）}
\begin{itemize}
    \item 逻辑表达式：$Y = \overline{AB}$
    \item 运算规则：有0出1，全1出0
\end{itemize}

\subsubsection{或非运算（NOR）}
\begin{itemize}
    \item 逻辑表达式：$Y = \overline{A+B}$
    \item 运算规则：全0出1，有1出0
\end{itemize}

\subsubsection{与或非运算（AOI）}
\begin{itemize}
    \item 逻辑表达式：$Y = \overline{AB+CD}$
    \item 运算规则：先两组与运算，再或运算，整体结果取反
\end{itemize}

\subsubsection{异或运算（XOR）}
\begin{itemize}
    \item 逻辑表达式：$Y = A \oplus B = A\overline{B}+\overline{A}B$
    \item 运算规则：输入相异输出1，输入相同输出0
\end{itemize}

\subsubsection{同或运算（XNOR）}
\begin{itemize}
    \item 逻辑表达式：$Y = A \odot B = AB+\overline{A}\overline{B}$
    \item 运算规则：输入相同输出1，输入相异输出0；同或是异或的反运算 $A\odot B=\overline{A\oplus B}$
\end{itemize}

\subsection{复合逻辑真值表（二输入）}
并排展示与非、或非、异或真值表，直观对比输入输出关系。
\begin{figure}[!htbp]
    \centering
    \begin{minipage}{0.32\textwidth}
        \centering
        \captionof{table}{与非真值表}
        \begin{tabular}{|c|c|c|}
            \hline
            $A$ & $B$ & $Y$ \\
            \hline
            0&0&1\\0&1&1\\1&0&1\\1&1&0\\
            \hline
        \end{tabular}
    \end{minipage}
    \hfill
    \begin{minipage}{0.32\textwidth}
        \centering
        \captionof{table}{或非真值表}
        \begin{tabular}{|c|c|c|}
            \hline
            $A$ & $B$ & $Y$ \\
            \hline
            0&0&1\\0&1&0\\1&0&0\\1&1&0\\
            \hline
        \end{tabular}
    \end{minipage}
    \hfill
    \begin{minipage}{0.32\textwidth}
        \centering
        \captionof{table}{异或真值表}
        \begin{tabular}{|c|c|c|}
            \hline
            $A$ & $B$ & $Y$ \\
            \hline
            0&0&0\\0&1&1\\1&0&1\\1&1&0\\
            \hline
        \end{tabular}
    \end{minipage}
\end{figure}

\subsection{复合逻辑门国标矩形符号}
独立单图，缩小缩放比例，横向不溢出版心；与或非门缩短横向偏移，紧凑排版。
\begin{figure}[!htbp]
    \centering
    \begin{tikzpicture}[scale=0.48]
        % 1.与非门
        \draw (0,0) rectangle (2,2);
        \node at (1,1) {\&};
        \draw (-0.8,1.5)--(0,1.5);
        \draw (-0.8,0.5)--(0,0.5);
        \draw (2,1)--(2.15,1);
        \draw (2.15,1) circle(0.12);
        \draw (2.27,1)--(2.8,1);
        \node[left] at (-0.8,1.5) {$A$};
        \node[left] at (-0.8,0.5) {$B$};
        \node[right] at (2.8,1) {$Y$};
        \node at (1,-0.8) {与非门};

        % 2.或非门
        \begin{scope}[xshift=4.8cm]
        \draw (0,0) rectangle (2,2);
        \node at (1,1) {$\geq1$};
        \draw (-0.8,1.5)--(0,1.5);
        \draw (-0.8,0.5)--(0,0.5);
        \draw (2,1)--(2.15,1);
        \draw (2.15,1) circle(0.12);
        \draw (2.27,1)--(2.8,1);
        \node[left] at (-0.8,1.5) {$A$};
        \node[left] at (-0.8,0.5) {$B$};
        \node[right] at (2.8,1) {$Y$};
        \node at (1,-0.8) {或非门};
        \end{scope}

        % 3.异或门
        \begin{scope}[xshift=9.6cm]
        \draw (0,0) rectangle (2,2);
        \node at (1,1) {$=1$};
        \draw (-0.8,1.5)--(0,1.5);
        \draw (-0.8,0.5)--(0,0.5);
        \draw (2,1)--(2.8,1);
        \node[left] at (-0.8,1.5) {$A$};
        \node[left] at (-0.8,0.5) {$B$};
        \node[right] at (2.8,1) {$Y$};
        \node at (1,-0.8) {异或门};
        \end{scope}

        % 4.同或门
        \begin{scope}[xshift=14.4cm]
        \draw (0,0) rectangle (2,2);
        \node at (1,1) {$=1$};
        \draw (-0.8,1.5)--(0,1.5);
        \draw (-0.8,0.5)--(0,0.5);
        \draw (2,1)--(2.15,1);
        \draw (2.15,1) circle(0.12);
        \draw (2.27,1)--(2.8,1);
        \node[left] at (-0.8,1.5) {$A$};
        \node[left] at (-0.8,0.5) {$B$};
        \node[right] at (2.8,1) {$Y$};
        \node at (1,-0.8) {同或门};
        \end{scope}

        % 5.国标与或非门（压缩横向宽度）
        \begin{scope}[xshift=19.2cm]
        \draw (0,1.8) rectangle (1.2,3.0);
        \node at (0.6,2.4) {\&};
        \draw (-0.6,2.7)--(0,2.7);
        \draw (-0.6,2.1)--(0,2.1);
        \node[left] at (-0.6,2.7) {$A$};
        \node[left] at (-0.6,2.1) {$B$};

        \draw (0,0) rectangle (1.2,1.2);
        \node at (0.6,0.6) {\&};
        \draw (-0.6,0.9)--(0,0.9);
        \draw (-0.6,0.3)--(0,0.3);
        \node[left] at (-0.6,0.9) {$C$};
        \node[left] at (-0.6,0.3) {$D$};

        \draw (2.0,0.9) rectangle (3.2,2.1);
        \node at (2.6,1.5) {$\geq1$};
        \draw (1.2,2.4)--(2.0,1.8);
        \draw (1.2,0.6)--(2.0,1.2);
        \draw (3.2,1.5)--(3.35,1.5);
        \draw (3.47,1.5) circle (0.12);
        \draw (3.59,1.5)--(4.0,1.5);
        \node[right] at (4.0,1.5) {$Y$};
        \node at (1.6,-0.8) {与或非门};
        \end{scope}
        \node[font=\bfseries] at (10,-2) {（a）国标矩形符号};
    \end{tikzpicture}
\caption{复合逻辑门国标矩形符号}
\label{fig:complex_gate_gb}
\end{figure}

\subsection{复合逻辑门ANSI异形符号}
独立第二张图，单独存放ANSI画法，避免单图内容过载；整体缩放适配页面。
\begin{figure}[!htbp]
    \centering
    \begin{circuitikz}[scale=0.48]
        % 与非门 NAND
        \draw (0,0) node[nand port] (nand1) {}
        (nand1.in 1)--++(-0.8,0)
        (nand1.in 2)--++(-0.8,0)
        (nand1.out)--++(0.8,0);
        \node at (0,-1.4) {与非门};

        % 或非门 NOR
        \draw (4.2,0) node[nor port] (nor1) {}
        (nor1.in 1)--++(-0.8,0)
        (nor1.in 2)--++(-0.8,0)
        (nor1.out)--++(0.8,0);
        \node at (4.2,-1.4) {或非门};

        % 异或门 XOR
        \draw (8.4,0) node[xor port] (xor1) {}
        (xor1.in 1)--++(-0.8,0)
        (xor1.in 2)--++(-0.8,0)
        (xor1.out)--++(0.8,0);
        \node at (8.4,-1.4) {异或门};

        % 同或门 XNOR
        \draw (12.6,0) node[xnor port] (xnor1) {}
        (xnor1.in 1)--++(-0.8,0)
        (xnor1.in 2)--++(-0.8,0)
        (xnor1.out)--++(0.8,0);
        \node at (12.6,-1.4) {同或门};

        % ANSI与或非门
        \draw (16.8,1) node[and port] (and1) {};
        \draw (16.8,-1) node[and port] (and2) {};
        \draw (19.5,0) node[or port] (or1) {};
        \draw (and1.out)--(18.2,0.5)--(or1.in 1);
        \draw (and2.out)--(18.2,-0.5)--(or1.in 2);
        \draw (or1.out)--(20.6,0);
        \filldraw[fill=white] (20.72,0) circle(0.12);
        \draw (20.84,0)--(21.4,0);
        \node at (19.5,-2) {与或非门};

        \node[font=\bfseries] at (10,-2.8) {（b）ANSI异形符号};
    \end{circuitikz}
\caption{复合逻辑门ANSI异形符号}
\label{fig:complex_gate_ansi}
\end{figure}
% ====================== 第四节优化结束 ======================

\section{逻辑代数的基本公式和常用公式}
\subsection{基本公式（布尔恒等式）}
\begin{table}[!htbp]
    \centering
    \caption{逻辑代数基本公式}
    \begin{tabular}{|c|c|c|c|}
        \hline
        序号 & 公式 & 序号 & 公式 \\
        \hline
        1  & $0 \cdot A = 0$        & 10 & $\overline{1}=0,\ \overline{0}=1$ \\
        2  & $1 \cdot A = A$        & 11 & $1 + A = 1$ \\
        3  & $A \cdot A = A$        & 12 & $0 + A = A$ \\
        4  & $A \cdot \overline{A} = 0$ & 13 & $A + A = A$ \\
        5  & $AB=BA$                & 14 & $A + \overline{A} = 1$ \\
        6  & $A(BC)=(AB)C$          & 15 & $A+B=B+A$ \\
        7  & $A(B+C)=AB+AC$         & 16 & $A+(B+C)=(A+B)+C$ \\
        8  & $\overline{AB}=\overline{A}+\overline{B}$ & 17 & $A+BC=(A+B)(A+C)$ \\
        9  & $\overline{\overline{A}}=A$ & 18 & $\overline{A+B}=\overline{A}\;\overline{B}$ \\
        \hline
    \end{tabular}
\end{table}

\subsubsection{核心运算法则}
\begin{itemize}
    \item 交换律：$AB=BA,\ A+B=B+A$
    \item 结合律：$(AB)C=A(BC),\ (A+B)+C=A+(B+C)$
    \item 分配律：$A(B+C)=AB+AC,\ A+BC=(A+B)(A+C)$
    \item 德·摩根定律：$\overline{AB}=\overline{A}+\overline{B},\ \overline{A+B}=\overline{A}\overline{B}$
\end{itemize}

\begin{table}[!htbp]
    \centering
    \caption{德摩根定律真值验证}
    \begin{tabular}{|c|c|c|c|}
        \hline
        $A$ & $B$ & $\overline{A+B}$ & $\overline{AB}$ \\
        \hline
        0&0&1&1\\0&1&0&1\\1&0&0&1\\1&1&0&0\\
        \hline
    \end{tabular}
\end{table}

\subsection{常用化简公式}
\begin{table}[!htbp]
    \centering
    \caption{逻辑化简常用公式}
    \begin{tabular}{|c|l|}
        \hline
        1 & $A + AB = A$ \\
        2 & $A + \overline{A}B = A + B$ \\
        3 & $AB + A\overline{B} = A$ \\
        4 & $A(A+B) = A$ \\
        5 & $AB+\overline{A}C+BC=AB+\overline{A}C$ \\
          & $AB+\overline{A}C+BCD=AB+\overline{A}C$ \\
        6 & $A\cdot\overline{AB}=A\overline{B},\ \overline{A}\cdot\overline{AB}=\overline{A}$ \\
        \hline
    \end{tabular}
\end{table}

\section{逻辑代数三大重要定理}
\subsection{代入定理}
\textbf{定理内容}：任何一个含有变量$A$的逻辑等式，若将等式中所有$A$均替换为任意逻辑函数$Y$，等式仍然成立。
\[
\text{若 } F(A,B,C,\dots)=G(A,B,C,\dots) \text{，则 } F(Y,B,C,\dots)=G(Y,B,C,\dots)
\]
\textbf{作用}：将两变量公式推广至多变量，拓展所有恒等式适用范围。

\subsection{反演定理（多变量推广）}
\[
\overline{A_1A_2\cdots A_n} = \overline{A_1}+\overline{A_2}+\cdots+\overline{A_n},\quad
\overline{A_1+A_2+\cdots+A_n} = \overline{A_1}\overline{A_2}\cdots\overline{A_n}
\]
求反函数$\overline{Y}$变换规则：$\cdot\leftrightarrow +$，$0\leftrightarrow1$，原变量$\leftrightarrow$反变量。

\subsection{对偶定理}
等式两边同时交换$\cdot\leftrightarrow +$、$0\leftrightarrow1$，变量不变，得到对偶式$Y^D$；原式成立则对偶式必然成立。

\section{逻辑函数及其描述方法}
\subsection{逻辑函数定义}
设输入逻辑变量为$A,B,C,\dots$，输出逻辑变量为$Y$，若输入与输出之间满足确定逻辑运算关系，则称$Y=F(A,B,C,\dots)$为逻辑函数；所有变量仅取0、1两种逻辑值。

\subsection{六大标准描述方式}
逻辑函数共有六类主流描述手段，各类形式可相互转换：
\begin{enumerate}
    \item \textbf{真值表}：表格形式，完整列出输入变量全部取值组合与对应输出值，直观无歧义；
    \item \textbf{逻辑代数式}：由与、或、非等逻辑运算组合构成的数学表达式，便于代数化简与推导；
    \item \textbf{逻辑图}：使用国标/ANSI标准逻辑门符号，按运算连接得到的电路图形，直接对应硬件电路；
    \item \textbf{波形图（时序图）}：随时间变化的高低电平波形，反映输入、输出动态时序关系；
    \item \textbf{卡诺图}：最小项几何排布方格图，专门用于快速化简逻辑函数；
    \item \textbf{硬件描述语言}：计算机专用文本描述方式，用于数字电路自动化设计。
\end{enumerate}

\subsection{各类描述方法的相互转换规则}
\subsubsection{真值表 $\leftrightarrow$ 逻辑代数式}
1. 真值表转标准与或式
\begin{itemize}
    \item 筛选表格中输出$Y=1$的全部输入变量组合；
    \item 每组组合写乘积项：变量取值为1写原变量，取值为0写反变量；
    \item 将全部乘积项相加，得到逻辑函数标准与或表达式。
\end{itemize}
2. 逻辑代数式转真值表
\begin{itemize}
    \item 列出输入变量全部$2^n$种取值组合；
    \item 将每组输入代入表达式计算输出$Y$，填入表格对应行。
\end{itemize}

\subsubsection{逻辑代数式 $\leftrightarrow$ 逻辑图}
1. 代数式绘制逻辑图：遵循运算优先级（非$\to$与$\to$或），从输入到输出逐级对应绘制非门、与门、或门、复合门；
2. 逻辑图书写代数式：从输入向输出逐层写出每一级门电路运算，合并整理得到完整表达式。

\subsubsection{真值表 $\leftrightarrow$ 波形图}
依据时序波形提取每一时刻输入电平组合，对照真值表得到对应输出电平，完成二者互相转换。

\subsubsection{真值表 $\leftrightarrow$ 卡诺图}
将真值表中每组输入组合对应的输出值，填入卡诺图对应最小项方格内，二者完全等价。

\subsection{逻辑函数标准形式}
逻辑函数存在两类标准表达式，以最小项、最大项为基础构成。
\subsubsection{最小项及标准与或式}
\begin{itemize}
    \item \textbf{定义}：$n$变量逻辑函数中，最小项$m_i$是包含全部$n$个变量的乘积项，每个变量以原变量或反变量形式出现且仅出现一次。$n$个变量共有$2^n$个最小项，按变量二进制取值换算十进制得到编号$i$。
    \item \textbf{核心性质}：仅一组变量取值能让对应最小项输出为1；全体最小项之和恒为1；任意两个不同最小项相乘结果为0；逻辑相邻的最小项可以合并消去一对互补变量。
    \item \textbf{标准形式}：将所有输出$Y=1$对应的最小项相加，写成 $Y=\sum m_i$ 的最小项之和形式。
\end{itemize}

\subsubsection{最大项及标准或与式}
\begin{itemize}
    \item \textbf{定义}：$n$变量逻辑函数中，最大项$M_i$是包含全部$n$个变量的或项，每个变量以原变量或反变量形式出现且仅出现一次。$n$个变量共有$2^n$个最大项，按变量使最大项取0的二进制取值换算十进制得到编号$i$。
    \item \textbf{核心性质}：仅一组变量取值能让对应最大项输出为0；全体最大项之积恒为0；任意两个不同最大项相加结果为1；逻辑相邻的最大项可以合并消去一对互补变量。
    \item \textbf{标准形式}：将所有输出$Y=0$对应的最大项相乘，写成 $Y=\prod M_i$ 的最大项之积形式。
\end{itemize}

\subsubsection{最小项与最大项的关系}
相同下标编号的最小项和最大项互为反量：$M_i=\overline{m_i}$，两种标准形式可以互相转换。

\section{逻辑函数化简}
\subsection{公式化简法}
利用逻辑代数基础公式，通过合并项、消去冗余项、配项法消去多余变量与冗余项。

\subsection{卡诺图化简法}
\subsubsection{卡诺图结构与逻辑相邻规则}
卡诺图将$n$变量最小项按几何相邻等价逻辑相邻排布，支持2/3/4/5变量，可由低变量卡诺图折叠生成。
逻辑相邻包含四类场景：左右方格相邻、上下方格相邻、左右边界循环相邻、四角方格相邻；仅逻辑相邻的方格才能合并圈组，不相邻方格禁止圈在一起。

\subsubsection{圈组化简核心原则}
1. 圈内方格数量必须为$2^k$：$1,2,4,8,16\cdots$；
2. 圈尽可能大，单个圈包含1越多，消去变量数量越多；
3. 圈的总数尽可能少，最终与或式项数最少；
4. 所有标1的方格至少被圈一次，允许方格被多个圈重复覆盖；
5. 圈组完成后，每个圈保留圈中不变的变量，相加得到最简与或式；
6. 化简结果不唯一，不同合法圈法可得到不同最简表达式。

\subsubsection{圈0化简技巧}
若卡诺图中0数量远少于1，可直接合并所有0得到反函数$Y'$，再对$Y'$整体取反得到原函数最简式，大幅简化运算。

\subsection{含无关项的卡诺图化简}
\subsubsection{无关项分类定义}
\begin{itemize}
    \item \textbf{约束项}：输入变量受物理场景限制，永远不会出现的取值组合，对应最小项恒等于0；
    \item \textbf{任意项}：输入该组合时电路输出无功能要求，输出可取0或1，不影响整体逻辑功能；
    \item \textbf{无关项统称}：约束项、任意项合并记为$\sum d$，逻辑函数标准写法：$Y=\sum m + \sum d$，约束条件$\sum d=0$。
\end{itemize}

\subsubsection{无关项化简规则}
1. 卡诺图中无关项方格标记为$\times$；
2. 化简时可按需将$\times$当作1或0处理；
3. 取舍目标：使可圈出的矩形圈更大，最终表达式最简。

\subsection{多输出逻辑函数化简}
多输出电路包含多个独立输出$Y_1,Y_2,Y_3\dots$，化简核心目标不是单独每个输出最简，而是寻找多个输出之间的公共与项，共享中间门电路，让整体电路总门、总输入引脚数量最小，降低硬件成本。

\subsection{逻辑函数形式变换}
数字工程中常统一采用单一类型逻辑门搭建电路，需利用反演律对标准与或式做形式转换：
\begin{enumerate}
    \item 与或式 $\rightarrow$ 与非-与非式：整体两次取反，内层用德摩根拆分；
    \item 与或式 $\rightarrow$ 或非-或非式：先求反函数或与式，再整体两次取反；
    \item 可结合异或、同或运算进一步压缩表达式项数。
\end{enumerate}